\section*{Resumen}
\addcontentsline{toc}{section}{Resumen}
\justifying
\setlength{\parindent}{1.5em}


En la División de Ciencias y Artes para el Diseño (CyAD) de la Universidad Autónoma Metropolitana Unidad Azcapotzalco, los profesores entregan cada trimestre tres tipos de documentos: las \emph{cartas temáticas} que describen el contenido de cada curso, los \emph{requisitos de recuperación} para estudiantes en situación irregular y una \emph{autoevaluación} única por trimestre de su propia labor docente. Hasta ahora este proceso se realizaba con formularios de Google y hojas de cálculo, lo que implicaba trabajo repetitivo, dificultades para dar seguimiento y una alta probabilidad de errores, pues muchos de los datos eran ingresados manualmente.

Para atender esta necesidad se desarrolló una aplicación web dividida en dos partes que se comunican entre sí: una parte visible en el navegador (frontend), con la que interactúan las personas usuarias, y una parte oculta que administra la información en una base de datos (backend). La aplicación se organiza alrededor de cinco funciones principales: acceso seguro con diferentes tipos de usuario (profesor, administrador y público general); captura y publicación de cartas temáticas; captura y publicación de requisitos de recuperación; llenado de autoevaluaciones con cálculo automático del nivel de desempeño; y generación de reportes de cumplimiento para la División.

Además de lo anterior se agregaron dos elementos que enriquecen la propuesta original: una vista pública que permite consultar los documentos sin necesidad de iniciar sesión, filtrando por trimestre, programa y unidad de enseñanza-aprendizaje, y si no hay documentos disponibles, se muestra el siguiente mensaje: "Solo se muestran UEA con al menos un documento publicado. Si una UEA no aparece es porque ningún profesor ha publicado su documento para este periodo"; y un panel de catálogos donde el personal administrativo mantiene actualizados los departamentos, licenciaturas, posgrados, áreas y unidades de enseñanza-aprendizaje (UEAs), esta última con posibilidad de importar información desde archivos CSV con la estructura adecuada.

En este reporte se describe el proceso de desarrollo, los resultados obtenidos, las diferencias entre lo comprometido y lo entregado --- en particular la decisión de generar reportes en formato Excel en lugar de PDF --- y las mejoras propuestas para versiones futuras del sistema.

\newpage
